\section{Workflow thực chiến}

\subsection{Cài đặt môi trường}

Repo hiện có hai cụm phụ thuộc chính.

\subsubsection{Đường XLA}

\begin{lstlisting}[style=rae]
conda create -n rae python=3.10 -y
conda activate rae
pip install uv
uv pip install torch~=2.5.0 torch_xla[tpu]~=2.5.0 torchvision==0.20.1 -f https://storage.googleapis.com/libtpu-releases/index.html
uv pip install timm==0.9.16 accelerate==0.23.0 torchdiffeq==0.2.5 wandb scipy torch-fidelity
uv pip install "numpy<2" "transformers==4.57.1" einops
\end{lstlisting}

\subsubsection{Phần bổ sung cho JAX / NNX}

\begin{lstlisting}[style=rae]
uv pip install "jax[cuda12]==0.5.1" flax==0.10.4 optax==0.2.4 orbax-checkpoint==0.11.16
uv pip install ml-collections clu absl-py etils huggingface_hub
\end{lstlisting}

Lần chạy JAX đầu tiên sẽ tự bootstrap backend \texttt{diffuse\_nnx} vào thư mục
\texttt{\~{}/.cache/rae\_jax/diffuse\_nnx}.
Các notebook JAX trên Kaggle không còn lặp lại từng lệnh \texttt{uv pip install}
này; chúng đặt \texttt{UV\_PROJECT\_ENVIRONMENT=/tmp/.venv},
\texttt{UV\_CACHE\_DIR=/tmp/uv-cache}, rồi chạy \texttt{uv sync -q} theo
\texttt{pyproject.toml} của repo trước các cell phụ thuộc package.
Repo này tự vá backend \texttt{diffuse\_nnx} để import các model Dinov2 từ
subpackage của \texttt{transformers}, nên nên dùng
\texttt{transformers==4.57.1} cho đường này.
Patch này cũng chuyển \texttt{google-cloud-storage} sang import lười, nên đường
Stage 1 với RAE/DINO không cần cài package đó trừ khi bạn thực sự dùng các
thành phần backend có tải asset từ GCS.
Patch bootstrap này cũng sửa khởi tạo EMA của backend để EMA bắt đầu từ bản
copy của model hiện thời thay vì cây tham số toàn số 0. Nếu bạn resume một
checkpoint Orbax cũ được tạo trước khi có fix này, checkpoint đó vẫn mang EMA
đã lưu sẵn trước đó.
Adapter cũng tự suy ra \texttt{downsample\_factor} và \texttt{latent\_channels}
từ \texttt{misc.latent\_size}, nên bước preview sample và FID của backend vẫn
chạy ở latent resolution thay vì lỡ cấp phát noise theo image resolution.

\subsection{Chuẩn bị model và dữ liệu}

\subsubsection{Model}

\begin{lstlisting}[style=rae]
hf download nyu-visionx/RAE-collections --local-dir models
\end{lstlisting}

\subsubsection{Dữ liệu}

Hầu hết các script dùng dữ liệu dạng \texttt{ImageFolder}. Ví dụ:

\begin{lstlisting}[style=rae]
dataset_root/
  class_a/
    0001.png
  class_b/
    0002.png
\end{lstlisting}

\subsection{Kiểm tra Stage 1 bằng reconstruction}

\subsubsection{Đường XLA / PyTorch}

\begin{lstlisting}[style=rae]
python src/stage1_sample.py \
  --config <config> \
  --image assets/pixabay_cat.png \
  --output recon.png
\end{lstlisting}

\subsubsection{Đường JAX}

\begin{lstlisting}[style=rae]
python3 src_jax/stage1_sample.py \
  --config <config> \
  --image assets/pixabay_cat.png \
  --output recon_jax.png
\end{lstlisting}

Lệnh này phù hợp để xác nhận:

\begin{itemize}[leftmargin=1.5em]
  \item checkpoint decoder Stage 1 có load được hay không,
  \item latent shape có đúng hay không,
  \item preprocessing của encoder có khớp hay không.
\end{itemize}

\subsubsection{Tính latent stat cho Stage 1 trên đường JAX}

Với dataset mới như CelebA, nên tạo một file bootstrap identity stat trước
(\texttt{mean = 0}, \texttt{var = 1}), rồi chạy:

\begin{lstlisting}[style=rae]
python3 src_jax/build_stage1_stats.py \
  --config <stage1_config> \
  --input <train_imagefolder> \
  --output <stage1_stat.pt> \
  --batch-size 16 \
  --set stage_1.params.normalization_stat_path=<bootstrap_identity_stat.pt>
\end{lstlisting}

File \texttt{stage1\_stat.pt} tạo ra vẫn cùng định dạng với repo gốc, nên dùng
lại được cho cả đường \texttt{src/} lẫn \texttt{src\_jax/}.
Các lệnh Stage 1 thuần ở JAX có thể chạy với YAML chỉ khai báo
\texttt{stage\_1}; chúng không bắt buộc có \texttt{stage\_2.target}.

\subsubsection{Reconstruction cả thư mục trên đường JAX}

\begin{lstlisting}[style=rae]
python3 src_jax/reconstruct_folder.py \
  --config <stage1_config> \
  --input <val_imagefolder> \
  --output-dir recon_jax_dir \
  --batch-size 8
\end{lstlisting}

Lệnh này hữu ích khi muốn export reconstruction set trước khi build FID reference
hoặc khi cần so sánh decoder Stage 1 giữa các checkpoint.

\subsection{Huấn luyện Stage 2 trên TPU bằng TorchXLA}

\begin{lstlisting}[style=rae]
python src/train.py \
  --config configs/stage2/training/ImageNet256/SiTDH-XL_DINOv2-B.yaml \
  --data-path <imagenet_train_split> \
  --results-dir results \
  --image-size 256 \
  --precision bf16
\end{lstlisting}

Trong mỗi optimizer step, train loop thực hiện:

\begin{enumerate}[leftmargin=1.8em]
  \item center crop và đưa ảnh thành tensor;
  \item mã hóa ảnh bằng RAE;
  \item tính transport loss trong latent space;
  \item step optimizer và scheduler trên TPU;
  \item update EMA;
  \item chạy logging, checkpointing, preview sample, validation và FID theo cadence.
\end{enumerate}

\subsection{Huấn luyện Stage 2 bằng JAX / NNX}

\begin{lstlisting}[style=rae]
python3 src_jax/train.py \
  --config configs/stage2/training/ImageNet256/SiTDH-XL_DINOv2-B.yaml \
  --data-path <imagenet_train_root> \
  --results-dir results_jax \
  --precision bf16 \
  --wandb
\end{lstlisting}

Các config ImageNet \texttt{SiTDH} được commit sẵn trên branch này giờ mặc định
đặt \texttt{training.num\_workers=16}, \texttt{training.prefetch\_factor=4},
\texttt{training.log\_rae\_latent\_stats=true}, và
\texttt{training.log\_activation\_stats=true}. Nếu bật block \texttt{eval},
adapter JAX cũng cho \texttt{eval.num\_workers} và
\texttt{eval.prefetch\_factor} kế thừa cùng bộ mặc định \texttt{16 / 4} trừ khi
override riêng.

Trên Kaggle TPU, nên bắt đầu từ bộ mặc định này nếu host chịu được; chỉ khi
runtime đó không ổn định mới hạ worker/prefetch bằng CLI. Adapter cũng tự tắt
backend TensorBoard summary writer trên Kaggle và giữ metric ở stdout cùng
wandb.

Điểm quan trọng:

\begin{itemize}[leftmargin=1.5em]
  \item vẫn dùng cùng file YAML của repo;
  \item hỗ trợ \texttt{--set key=value} để override nhanh từ CLI;
  \item hỗ trợ \texttt{--wandb-run-id <existing\_run\_id>} để bind một
  legacy workdir với đúng W\&B run cũ đúng một lần;
  \item nếu \texttt{stage\_2.ckpt} là file PyTorch \texttt{.pt}, adapter sẽ port
  trọng số sang NNX để khởi tạo;
  \item nếu \texttt{stage\_2.ckpt} là thư mục Orbax, adapter sẽ restore run JAX
  trước đó.
\end{itemize}

\subsection{Wandb logging}

Biến môi trường khuyến nghị:

\begin{lstlisting}[style=rae]
export ENTITY=<wandb_entity>
export PROJECT=<wandb_project>
export WANDB_KEY=<wandb_api_key>
\end{lstlisting}

Đường JAX cũng chấp nhận các tên \texttt{WANDB\_ENTITY} và
\texttt{WANDB\_API\_KEY}. Adapter sẽ bridge các tên legacy ở trên nếu chúng đã
được export sẵn. Với \texttt{src\_jax/train.py}, mỗi workdir mới giờ sẽ ghi
\texttt{wandb\_run.json} để khóa đúng W\&B run id của workdir đó. Run mới dùng
\texttt{resume="never"}; các lần resume sau sẽ đọc lại file này, giữ đúng run
id đã bind, rồi tự rewind W\&B history về latest \texttt{checkpoint\_*} step
qua \texttt{resume\_from}. Nhờ đó phần metric đã log sau checkpoint cuối của
session cũ sẽ không bị giữ lại sai lịch sử. Nếu resume một Orbax workdir cũ chưa có
\texttt{wandb\_run.json}, cần truyền \texttt{--wandb-run-id <existing\_run\_id>}
một lần để bind đúng historical run trước khi train tiếp.

\subsection{FID trong train loop}

Đường XLA hỗ trợ online FID trong \texttt{src/train.py}. Ví dụ:

\begin{lstlisting}[style=rae]
eval:
  fid_ref: /path/to/reference_stats.npz
  fid_every: 25000
  fid_num_samples: 4096
  fid_per_proc_batch_size: 4
  fid_batch_size: 128
  fid_device: cpu
  fid_num_threads: 96
\end{lstlisting}

Đường JAX giờ cũng dùng block \texttt{eval} để chạy validation loss trên held-out
\texttt{ImageFolder}, log \texttt{eval/ema\_loss} mặc định và thêm
\texttt{eval/model\_loss} nếu bật \texttt{eval\_model}. Với online FID, nên xây
\texttt{fid\_ref} bằng detector backend-native qua:

\begin{lstlisting}[style=rae]
python3 src_jax/build_fid_stats.py \
  --input /path/to/reference_images \
  --output /path/to/reference_stats.pkl \
  --batch-size 64 \
  --num-workers 8
\end{lstlisting}

Trên đường JAX, metric \texttt{FID-4K (cfg=...)} mặc định vẫn chấm EMA. Nếu bật
\texttt{fid\_eval\_model}, adapter sẽ log thêm metric riêng
\texttt{FID-4K/model (cfg=...)} cho online model để bạn đối chiếu trực tiếp với
panel \texttt{network\_samples}. Nếu bật hai cờ diagnostics trong block
\texttt{training}, đường JAX còn log thêm \texttt{train\_rae\_latent\_rms},
\texttt{train\_rae\_latent\_var}, \texttt{train\_sitdh\_output\_rms},
\texttt{train\_sitdh\_output\_var}, và RMS/variance cho từng block
\texttt{SiTDH} như \texttt{train\_sitdh\_act\_enc\_00\_rms} hoặc
\texttt{train\_sitdh\_act\_dec\_01\_var}.

\subsection{Sampling Stage 2}

\subsubsection{JAX một lệnh}

\begin{lstlisting}[style=rae]
python3 src_jax/sample.py \
  --config configs/stage2/sampling/ImageNet256/SiTDHXL-DINOv2-B_AG.yaml \
  --class-labels 207,360 \
  --output sample_jax.png
\end{lstlisting}

\subsubsection{JAX distributed}

\begin{lstlisting}[style=rae]
python3 src_jax/sample_ddp.py \
  --config configs/stage2/sampling/ImageNet256/SiTDHXL-DINOv2-B.yaml \
  --sample-dir samples_jax \
  --num-samples 50000 \
  --label-sampling equal \
  --fid-ref /path/to/reference_stats.npz
\end{lstlisting}

Các config sampling JAX được commit dưới dạng template. Trước khi chạy sampling,
cần tự điền \texttt{stage\_2.ckpt} và, nếu dùng autoguidance, cả
\texttt{guidance.guidance\_model.ckpt} bằng checkpoint \texttt{SiTDH} tương
thích.

\subsection{Đẩy artifact lên Hugging Face}

\subsubsection{Lệnh độc lập}

\begin{lstlisting}[style=rae]
python3 src_jax/push_hf.py \
  --path results_jax/<run_name> \
  --repo-id <hf_user_or_org>/<repo_name>
\end{lstlisting}

\subsubsection{Tích hợp ngay trong train}

\begin{lstlisting}[style=rae]
python3 src_jax/train.py \
  --config <config> \
  --data-path <train_root> \
  --hf-repo-id <hf_user_or_org>/<repo_name>
\end{lstlisting}

\subsection{Notebook Kaggle cho CelebA}

Repo hiện có notebook \texttt{raes-jax-celeba-kaggle.ipynb} để bám đúng flow
Kaggle:

\begin{itemize}[leftmargin=1.5em]
  \item clone repo và checkout branch \texttt{jax-sit-dh};
  \item đặt \texttt{UV\_PROJECT\_ENVIRONMENT=/tmp/.venv} và \texttt{UV\_CACHE\_DIR=/tmp/uv-cache};
  \item chạy \texttt{uv sync -q} từ root của repo;
  \item chạy các bước phụ thuộc package qua \texttt{uv run};
  \item chuyển CelebA thành \texttt{ImageFolder} \texttt{256x256};
  \item tạo bootstrap identity stat, build latent stat cho Stage 1;
  \item export reconstruction validation set, build \texttt{fid\_ref};
  \item ghi sẵn config CelebA cho \texttt{src\_jax/train.py}.
\end{itemize}

Nếu chạy trên Kaggle \texttt{TPU v5e-8}, dùng notebook
\texttt{raes-jax-celeba-kaggle-tpuv5e8-sitdh-s.ipynb}. Bản này cố định biến thể
CelebA Stage 2 là \texttt{SiTDH-S}, sync dependency của repo vào
\texttt{/tmp/.venv}, tự xóa cờ executable-stack mà Kaggle có thể chặn trên
\texttt{jaxlib}, chạy bước xác nhận TPU trong một tiến trình Python mới, giữ
Stage 1 và Stage 2 ở TPU, đồng thời build \texttt{fid\_ref} bằng đúng detector
backend dùng cho online FID, và giữ cadence checkpoint mặc định ở mốc
\texttt{210000} step.

Nếu muốn giữ nguyên flow đó nhưng đổi Stage 2 sang biến thể DH
\texttt{SiTDH-B},
dùng notebook \texttt{raes-jax-celeba-kaggle-tpuv5e8-sitdh-b.ipynb}. Notebook
này ghi config CelebA thành
\texttt{CelebA256\_SiTDH-B\_DINOv2-B\_jax\_tpuv5e8.yaml}, đặt backbone DH thành
\texttt{hidden\_size=[768, 2048]}, \texttt{depth=[12, 2]},
\texttt{num\_heads=[12, 16]}, bật \texttt{use\_pos\_embed}, tắt label dropout
bằng \texttt{class\_dropout\_prob=0.0}, đặt tên run/project mặc định theo biến
thể \texttt{SiTDH-B}, và
giữ cùng cadence checkpoint \texttt{210000} step.

Nếu đã có thư mục run Orbax của \texttt{SiTDH-B} và cần train tiếp, dùng
notebook \texttt{raes-jax-celeba-kaggle-tpuv5e8-sitdh-b-resume.ipynb}. Notebook
này đã được rút xuống thành bản resume-only cho trường hợp bắt đầu từ output
archive của notebook cũ: cell đầu chạy lệnh \texttt{unzip -o
/kaggle/input/notebooks/kieuhongquan/rae-jax/\_output\_.zip -d
/kaggle/working}, sau đó mới kiểm tra repo đã có sẵn, chạy \texttt{uv sync},
fix \texttt{jaxlib}, lấy secret wandb, kiểm tra path, rồi tự tìm thư mục
\texttt{CelebA256\_SiTDH-B\_DINOv2-B\_jax\_tpuv5e8-*} mới nhất bên dưới
\texttt{/kaggle/working/results\_jax\_tpu/}, rồi mới tìm thư mục
\texttt{checkpoint\_<step>/} mới nhất trong workdir đó trước khi truyền
\texttt{--workdir} vào \texttt{src\_jax/train.py}. Runtime tiếp tục restore
\texttt{latest\_step()} từ checkpoint mới nhất, nên notebook này không còn
hardcode cả thư mục run có timestamp lẫn \texttt{checkpoint\_100000}. Nếu
workdir đích là legacy run được tạo trước khi cơ chế
\texttt{wandb\_run.json} tồn tại, cần thêm
\texttt{--wandb-run-id <existing\_run\_id>} một lần để notebook resume đúng
chính run W\&B cũ đó.

Branch này cũng giữ song song workflow CelebA-HQ:
\texttt{raes-jax-celebahq-kaggle.ipynb},
\texttt{raes-jax-celebahq-kaggle-tpuv5e8-sitdh-s.ipynb},
\texttt{raes-jax-celebahq-kaggle-tpuv5e8-sitdh-b.ipynb}, và
\texttt{raes-jax-celebahq-kaggle-tpuv5e8-sitdh-b-resume.ipynb}. Chúng dùng lại
pipeline JAX/Kaggle tương tự nhưng export dữ liệu CelebA-HQ 256. Merge này
cũng khôi phục hai helper export riêng là
\texttt{src\_jax/export\_celebahq\_hf.py} và
\texttt{src\_jax/export\_celebahq\_tfds.py}.
